\section{Related Work}

% should we include this more general related work
\noindent \textbf{Novel-View Synthesis (NVS) and Neural Rendering.} 
NVS tackles the problem of generating accurate renderings of 3D scenes from unseen viewpoints given a set of training images and camera poses. 
NVS traces back to the Structure from Motion (SfM) works \cite{snavely2026phototourism, Hartley2004sfm, brown2005unsupervised} and Multiview Stereo (MVS) \cite{seitz2006comparison, goesele2007multi, huang2018deepmvs}, which served as the basis of several groundbreaking NVS methods \cite{DeepBlending2018, eisemann2008floating, chaurasia2013depth, KPLD21point}. 
However, these approaches generally require storing hundreds and thousands of input images for reprojection and blending, which leads to massive memory requirements and sometimes fails to reconstruct undersampled regions. 

Recent advances in neural rendering \cite{tewari2022advances} have made great strides in improving the quality of 3D reconstruction and novel-view rendering. 
Neural rendering algorithms can be roughly classified by the underlying 3D representation, ranging from point clouds~\cite{KPLD21point, aliev2019point, xu2022point}, voxels~\cite{yu_and_fridovichkeil2021plenoxels, hu2023multiscale, yu2021plenoctrees}, meshes~\cite{neumesh, kato2018meshrenderer, Lombardi2018deepappearance, hu2021worldsheet, guo2023vmesh}, to implicit representations using MLPs~\cite{mildenhall2020nerf, mueller2022instant, sitzmann2019srns,attal2022learning}. 
Recently, 3D Gaussian Splatting \cite{kerbl3Dgaussians} (3DGS) has emerged as the state-of-the-art NVS technique. 
3DGS achieves high image quality for NVS, is fast to optimize and render, and leverages an explicit Gaussian primitive representation, which makes the method suitable for a wide variety of tasks such as surface extraction \cite{guedon2023sugar, Huang2DGS2024}, human avatar manipulation \cite{qian20233dgsavatar, shao2024splattingavatar}, ambient dynamic modeling~\cite{shih2024modeling}, and object and scene editing \cite{guedon2024frosting, gaussian_grouping}. 

\begin{figure*}[!t]
    \centering
\includegraphics[width=\textwidth]{sec/figures/model_pipeline.pdf}
    \caption{\textbf{Textured Gaussians model pipeline.} 
    Our method consists of three major components: ray-Gaussian intersection, RGBA texture mapping, and a generalized Gaussian appearance model. 
    To render the color of a pixel $\mathbf{p}$, we first trace a ray from the camera center $\mathbf{o}$ to the pixel to intersect with 3D Gaussians in the scene. 
    Then, we query texture and alpha values, $\mathbf{c}_\text{tex}$ and $\alpha_\text{tex}$, from the per-Gaussian RGBA texture maps using the ray-Gaussian intersection point $\mathbf{x}$. 
    Finally, given the retrieved spatially varying color and alpha values, we alpha-composite the color and alpha values of Gaussians that are hit by the pixel ray using the generalized Gaussian appearance model. }
    \label{fig:model_pipeline}
\end{figure*}

Our method augments 3DGS with alpha and texture mapping commonly used in mesh-based 3D representations. 
During optimization and rendering, we intersect outgoing rays from pixels with Gaussians in the scene. 
Intersection points are then used to query per-Gaussian texture map values via UV mapping and compute RGB color and alpha blending values. 
This enables each Gaussian to represent complex textures and shapes, significantly improving NVS quality.

% weak ending

\noindent \textbf{Appearance Modeling in 3D Gaussian Splatting.}  
In 3DGS, scene geometry is represented by the position, rotation, and scale of 3D Gaussians. By adjusting these properties, 3DGS can represent intricate geometric structures. 
On the other hand, appearance is modeled using per-Gaussian opacity values and spherical harmonics coefficients. However, pixels within a projected Gaussian are always shaded with the same color up to a Gaussian falloff factor, greatly limiting the expressivity of individual Gaussians. Therefore, recent works have explored ways that allow Gaussians to represent spatially varying features. \citet{huang2024spatial} defined per-Gaussian SH coefficients for color and opacity and used ray-Gaussian intersections to determine the viewing directions within the Gaussian extent. This enables smooth variation of colors and opacities across pixels covered by a single Gaussian but prevents the reconstruction of higher-frequency details. \citet{xu2024texturegs} disentangled appearance and geometry in 3DGS for texture editing, but their method is limited to object-centric scenes with simple geometry due to the unit-sphere parameterization of textures.

Instead of using a global texture map, we enhance each 3D Gaussian with a local texture map on top of the SH coefficients. 
This allows for a more flexible texture optimization since individual Gaussians are not tied to a shared global texture map. 
Our method can, therefore, reconstruct objects with complex structures and real-world scenes. 
Furthermore, each Gaussian can also represent various shapes with the alpha channel in the texture map.

\noindent \textbf{Memory Efficient Gaussian Splatting. } 
In order to achieve even faster rendering time and compact storage requirements for potential applications on edge devices, there has been a plethora of recent work focusing on optimizing memory-efficient Gaussian Splatting models \cite{bagdasarian20243dgszipsurvey3dgaussian}. 
These algorithms either prune Gaussians and perform quantization of Gaussian attributes \cite{kerbl2024reducing, lee2024compact, navaneet2023compact3d, morgenstern2023compact, fan2023lightgaussian, Niedermayr_2024_CVPR} or exploit the structural relationships between Gaussians \cite{chen2024hash, tao2024scaffoldgs}. 

Our work is orthogonal and complementary to these works since we focus on improving the appearance modeling of Gaussians by redistributing the model size budget to the per-Gaussian texture maps given any optimized 3DGS model. 
These model compression algorithms can be easily integrated with the 3DGS pretraining stage in our optimization pipeline to achieve further gains in compactness.

\noindent \textbf{Concurrent Work.} Different from several concurrent works that have explored the idea of augmenting Gaussians with texture maps \cite{rong2024gstex, song2024hdgstextured2dgaussian, xu2024SuperGaussians, svitov2025billboardsplattingbbsplatlearnable}, we 1) additionally utilize alpha textures and 2) provide an in-depth analysis of various texture map variants (alpha-only, RGB, RGBA) that explores the balance between reconstruction capability and model size.
